\section{Correction 与迭代循环}

\subsection{Correction 环节}

当 Verifier 判定 solution 无效时，提取 bug report 送入 Correction 环节：
\begin{itemize}
    \item 输入：原始问题 + 当前 solution + bug report
    \item Prompt：``Below is a bug report. If you agree with it, fix the issues. Your output should follow the system prompt instructions.''
    \item 输出：修正后的新 solution
\end{itemize}

\subsection{完整迭代流程}

以第一道题为例：
\begin{enumerate}
    \item 初始探索：Step 1 + Step 2 生成初始解
    \item 验证：Verifier 判定 No（有 bug）
    \item 纠正：基于 bug report 生成新解 $\rightarrow$ 验证 No $\rightarrow$ 纠正 $\rightarrow$ 验证 No
    \item 第 4 次纠正后通过验证：correct count = 1
    \item 连续验证通过 5 次，最终接受
    \item 共经历 8 次迭代，外部循环仅 1 次 run
\end{enumerate}

\subsection{本章小结}

整个迭代循环展示了 Actor-Critic 模式的强大：即使初始解有多处错误，通过反复的验证-纠正循环，最终可以收敛到正确解。

